\chapterw{Lernziele: Rechnernetze}
\begin{todolist}
\item Ich kann den Unterschied zwischen Internet und Rechnernetz erklären und Beispiele für verschiedene Netzwerke nennen.
\item Ich kann die Funktionsweise von IP, IPv4 und IPv6 erklären und die Adressräume vergleichen.
\item Ich kann die Erreichbarkeit eines Hosts mit \texttt{ping} prüfen und typische Diagnoseschritte durchführen.
\item Ich kann die Funktionen und Arbeitsweisen von Routern und Switches in einem Netzwerk verstehen und erläutern.
\item Ich kann Subnetzmasken anwenden und einfache Subnetze planen.
\item Ich kann eine Routing-Tabelle interpretieren und für einfache Routing-Entscheidungen nutzen.
\item Ich kann eine Subnetzmaske erstellen und anwenden, um Netzwerke in verschiedene Subnetze zu unterteilen.
\item Ich kann die grundlegenden Konzepte von TCP beschreiben und dessen Einsatz im Kontext von Netzwerkprotokollen verstehen.
\item Ich kann die Rolle von Ports anhand einer Client-Server-Kommunikation erklären.
\item Ich kann Header und Payload eines Datenpakets beschreiben und deren Zweck erklären.
\item Ich kann grundlegende HTML-Strukturen erstellen und verstehen, wie HTML-Seiten über HTTP übertragen werden.
\item Ich kann den Ablauf eines Abrufs von Webinhalten über das Internet in einzelnen Schritten und den darin involvierten Protokollen erläutern.
\item Ich kann die Rolle eines DNS-Servers im Netzwerk erklären und wie DNS-Anfragen bearbeitet und beantwortet werden.
\item Ich kann die grundlegenden Schritte zur Fehlerbehebung einer nicht funktionierenden Verbindung in Filius durchführen und dokumentieren.
\item Ich kann eine lokale Chat-Verbindung über \texttt{localhost} aufbauen und testen.
\item Ich kann ein einfaches Socket-Programm in Python analysieren und anpassen (Client und Server).
\item Ich kann eine Chat-Verbindung im LAN einrichten und notwendige Einstellungen (Firewall/Ports) berücksichtigen.
 \item Ich kann die Funktionsweise eines Relay-Servers erklären und einfache Nachrichtenformate definieren.
 \item Ich kann mit Netzwerkanalyse-Tools (z.B. Wireshark) gängige Protokolle in einem Abruf identifizieren (DNS, TCP, HTTP).
 \item Ich kann Chancen und Risiken zentraler Serverarchitekturen hinsichtlich Datenschutz und Souveränität reflektieren.
 \item Ich kann konkrete Massnahmen zur Datenminimierung und zum Schutz der Privatsphäre in einem Chat-Design vorschlagen.
\end{todolist}
